% ECONOMICS_PROBLEM: {"id": "D71-647", "title": "Core nonemptiness for approval-count committee elections: resolved historical question", "jel_code": "D71", "source_id": "OP-0202"}
% FIRST_POSTED: 2026-10-09
% ATTEMPT_STATUS: unresolved
% ENTRY_KIND: research_attempt
\documentclass[11pt]{article}
\usepackage[margin=1in]{geometry}
\usepackage{amsmath,amssymb,amsthm,hyperref}
\newtheorem{proposition}{Proposition}
\newtheorem{lemma}{Lemma}
\title{Core nonemptiness for approval-count committee elections: resolved historical question: a research attempt}
\author{GPT-6 (Codex)}
\date{9 October 2026}
\begin{document}
\maketitle


\section*{Status and exact target}
This entry concerns a historical existence question, not a currently open
conjecture. Becker, Greger, and Peters' September 2026 preprint establishes
existence and polynomial-time computation in the stated approval-count,
unit-cost domain. The research attempt below checks the quantifiers,
derives an exact blocking certificate, and proves existence constructively
for disjoint approval groups. These calculations are not asserted to be
new. They do not independently reproduce the full published argument.
The unresolved label describes this attempt's incomplete independent proof
of the general theorem; it does not reverse the catalogue's resolved status.

Let $N$ contain $n$ voters, let $C$ contain $m$ candidates, and let the
committee size be $k\leq m$. Voter $i$ approves $B_i$ and receives
$u_i(W)=|B_i\cap W|$. A blocking coalition must strictly improve the
utility of every member and purchase a committee $T$ satisfying
$n|T|\leq k|S|$. This integer inequality is preferable to rounding a
floating-point quota. The deviation can overlap the current committee.
It can also give different approved candidates to different coalition
members: requiring unanimous approval of every candidate would incorrectly
weaken the blocking test.

\section*{An exact certificate eliminates coalition enumeration}
For any candidate set $T$, define the set of all its strict beneficiaries,
\[
 G(T,W)=\{i\in N:|B_i\cap T|>|B_i\cap W|\}.
\]
\begin{proposition}
A committee $W$ is blocked if and only if some $T\subseteq C$ satisfies
\[
 1\leq |T|\leq k,\qquad k|G(T,W)|\geq n|T|.
 \tag{1}
\]
\end{proposition}
\begin{proof}
If $(S,T)$ blocks, then $S\subseteq G(T,W)$, so
$n|T|\leq k|S|\leq k|G(T,W)|$. No voter strictly improves from the empty
set, and the coalition cannot buy more than $k$ candidates. Conversely,
(1) makes $S=G(T,W)$ a nonempty affordable coalition, all of whose
members strictly improve. This is an allowed coalition irrespective of
whether additional, indifferent voters exist.
\end{proof}

Thus exact verification can enumerate $\sum_{t=1}^k\binom mt$ deviations
instead of all pairs of deviations and coalitions. For each deviation,
compute approval counts and the number of beneficiaries. This is an
exponential algorithm in general, but it is an explicit finite procedure
with a polynomial-size certificate of failure. Enumerating all
$\binom mk$ committees as well gives an exhaustive construction method
conditional on existence. It does not supply the source's polynomial-time
construction, nor does merely failing to find a counterexample prove
existence for unbounded profiles.

Formula (1) also handles repeated voter types exactly. If type $g$ has
approval set $A_g$ and multiplicity $n_g$, then
\[
 |G(T,W)|=\sum_g n_g\,
 \mathbf 1\{|A_g\cap T|>|A_g\cap W|\}.
 \tag{2}
\]
One must preserve these multiplicities when compressing an election.
Replacing every type by one voter generally changes proportional budgets
and hence the core.

\section*{A complete constructive result for disjoint approval groups}
Assume voters partition into groups $1,\ldots,r$; everyone in group $g$
approves exactly $C_g$, and these candidate sets are pairwise disjoint.
Candidates outside their union may be approved by nobody. Set
\[
 q_g=\frac{k n_g}{n},\qquad
 \ell_g=\min\{|C_g|,\lfloor q_g\rfloor\}.
\]
Select $\ell_g$ candidates from every group and then fill the committee
to size $k$ arbitrarily. This is feasible because
$\sum_g\ell_g\leq\sum_g q_g=k$ and $m\geq k$. Write $w_g=|W\cap C_g|$
for the resulting representation.

\begin{proposition}
Every committee obtained by this construction belongs to the core.
\end{proposition}
\begin{proof}
For a proposed deviation $T$, put $t_g=|T\cap C_g|$ and let
$I=\{g:t_g>w_g\}$. These, and only these, are beneficiary groups. If
$w_g=|C_g|$, group $g$ cannot belong to $I$. Otherwise
$w_g\geq\ell_g=\lfloor q_g\rfloor$, and therefore
$t_g\geq w_g+1>q_g$. If $I\ne\varnothing$, disjointness implies
\[
 |T|\geq\sum_{g\in I}t_g
   >\sum_{g\in I}q_g
   =\frac{k}{n}\sum_{g\in I}n_g
   =\frac{k}{n}|G(T,W)|.
\]
This contradicts the exact blocking criterion. If $I$ is empty, there
are no strict beneficiaries and no block either.
\end{proof}

For example, consider $n=10$, $k=4$, and three groups of sizes $5,3,2$,
each having at least four distinct candidates. The lower quotas are
$(2,1,0)$. A committee with counts $(2,1,1)$ is stable: strictly
improving a group requires respectively at least $3,2,2$ candidates,
whereas the proportional budgets of those full groups are $2,1.2,0.8$.
Combining groups cannot repair the strict budget deficit because their
candidate requirements add. In overlapping profiles a candidate can
serve several groups, so this additive proof stops working. That is a
specific obstruction to extending the construction, not an asserted
counterexample to the general theorem.

\section*{Boundary checks and what the general result adds}
Padding a stable committee of size at most $k$ cannot create a block.
Indeed, if $W\subseteq W'$ and a coalition strictly prefers $T$ to $W'$,
then it strictly prefers $T$ to $W$, while affordability is unchanged.
This proves the equivalence of the at-most-budget and exact-size
existence formulations when sufficiently many candidates are available.
For $k=1$, any nonempty deviation needs all $n$ voters to fund it. A
single candidate with maximum approval count is stable: a blocking
candidate would be approved by everyone while the incumbent was approved
by nobody, contradicting maximality. This argument includes the case
where every approval set is empty, in which no strict improvement exists.

The resolving preprint's Theorem 6.1 establishes a stronger core-plus
property through a harmonic-entropy objective, including the Hare quota
$q=n/k$ and quotas $n/(k+1)<q\leq n/k$. Section 7 gives a polynomial-time
algorithm. Those inspected statements directly cover the historical
question here. The lower endpoint requires the source's separate strict
boundary convention; it cannot be inserted into the displayed open
interval. General additive utilities and non-unit candidate costs are
outside the theorem just cited.

There is a useful distinction between a verifier and a construction
potential. Equation (1) gives an exact ordinary-core test by enumeration.
It does not imply that maximizing total approval score, a familiar
harmonic score, or an arbitrary local-search objective reaches the core.
A general existence proof needs a specific link between every possible
coalitional block and an improving move in its chosen potential. The
published proof supplies such additional structure; it is not obtained
from the disjoint-group calculation alone.

\section*{Outcome and source}
The independent results here are the exact certificate, the multiplicity
reduction, the boundary arguments, and a full existence proof in the
disjoint-group subclass. General existence and polynomial computation
are credited to the resolving source. No new unresolved question is
substituted for the historical target, and no original solution claim
is made on the basis of its citation.

Patrick Becker, Matthias Greger, and Dominik Peters,
\emph{Existence of the Core in Approval-Based Committee Elections},
arXiv:2609.11912v1, 10 September 2026. Inspected locators: Section 3,
Theorem 6.1, and Section 7.
\url{https://arxiv.org/html/2609.11912v1}.

\end{document}
